\title{DeepSeekMath: Pushing the Limits of Mathematical Reasoning in Open Language Models}

\begin{abstract}
The inherently intricate and structured characteristics of mathematical reasoning make it particularly challenging for language models to master. This study presents DeepSeekMath~7B, a model developed by further pre-training DeepSeek-Coder-Base-v1.5 7B with 120 billion math-centric tokens derived from Common Crawl, combined with natural language and code corpora. Remarkably, DeepSeekMath~7B attains a 51.7\% score on the MATH benchmark at a competition level, without utilizing external tools or ensembling methods, thereby nearing the accuracy rates of Gemini-Ultra and GPT-4. When employing self-consistency with 64 samples, DeepSeekMath~7B reaches 60.9\% on MATH. The advancement in mathematical reasoning exhibited by DeepSeekMath~can be traced to two primary innovations: First, we leverage the extensive utility of open-access web resources through a carefully designed data curation process. Second, we propose Group Relative Policy Optimization (GRPO), an extension of Proximal Policy Optimization (PPO), which not only augments mathematical reasoning performance but also streamlines PPO's memory demands.
\end{abstract}

\section{Introduction}

Large language models (LLM) have transformed the landscape of mathematical reasoning within artificial intelligence, driving notable progress on tasks such as the quantitative reasoning benchmark \citep{MATH} as well as the geometry reasoning benchmark \citep{trinh2024solving}. In addition to their benchmarking success, these models have also served as valuable aids to humans working through intricate mathematical problems \citep{tao}. Despite these achievements, state-of-the-art systems like GPT-4 \citep{gpt4} and Gemini-Ultra \citep{gemini} remain closed to public access, and the current roster of open-source alternatives lags substantially in terms of effectiveness.

In this work, we present DeepSeekMath, a specialized language model designed for mathematical reasoning that notably surpasses existing open-source models and comes close to the academic benchmark scores achieved by GPT-4. Our approach involves constructing the DeepSeekMath Corpus—a comprehensive, high-quality pre-training dataset consisting of 120B math-specific tokens. This corpus is derived from the Common Crawl (CC) by employing a fastText-based classifier \citep{joulin2016fasttext}. Initially, the classifier is trained on positive examples sourced from OpenWebMath \citep{openwebmath}, complemented by a broad array of unrelated webpages as negative samples. The trained classifier is subsequently used to extract further positive examples from the CC, which are then curated through manual annotation. Incorporating these refined examples, the classifier is retrained to further enhance its efficacy. Our evaluation demonstrates the superior quality of this large-scale dataset: our primary model, DeepSeekMath-Base 7B, achieves a score of 64.2\% on GSM8K \citep{gsm8k} and 36.2\% on the advanced MATH benchmark \citep{MATH}, outperforming Minerva 540B \citep{minerva}. Furthermore, since the DeepSeekMath Corpus is inherently multilingual, we observe performance gains on Chinese mathematical tasks \citep{wei2023cmath,agieval}. We anticipate that our methodology for mathematical data curation serves as a valuable reference point for subsequent research, with substantial opportunities for further advancements.

DeepSeekMath-Base is built upon DeepSeek-Coder-Base-v1.5 7B \citep{deepseek-coder}, since leveraging a model pretrained on code proves to be more effective than initializing from a standard large language model. Additionally, our findings reveal that mathematical training leads to performance gains not just in mathematics, but also on the MMLU \citep{mmlu} and BBH benchmarks \citep{bbh}. This suggests that such training enhances both mathematical proficiency and general reasoning skills within the model.

Following the pre-training phase, we conduct mathematical instruction tuning on DeepSeekMath-Base utilizing datasets focused on chain-of-thought \citep{cot}, program-of-thought \citep{pot,pal}, and tool-integrated reasoning \citep{tora}. The resulting model, DeepSeekMath-Instruct 7B, outperforms all other models with a 7B parameter count and demonstrates performance on par with open-source instruction-tuned models at the 70B scale.

In addition, we present Group Relative Policy Optimization (GRPO), a novel reinforcement learning (RL) algorithm inspired by Proximal Policy Optimization (PPO) \citep{schulman2017proximal}. Unlike standard approaches that rely on a critic network, GRPO eliminates the critic, instead deriving baseline estimates from aggregated group scores, which leads to a marked decrease in the computational requirements for training. Utilizing only a fraction of English instruction tuning data, GRPO demonstrates notable performance gains over the robust DeepSeekMath-Instruct model, both on in-domain benchmarks (GSM8K: 82.9\% $\rightarrow$ 88.2\%, MATH: 46.8\% $\rightarrow$ 51.7\%) and on out-of-domain mathematical evaluations (for example, CMATH: 84.6\% $\rightarrow$ 88.8\%) during reinforcement learning. Furthermore, we propose a cohesive framework that contextualizes various approaches—including Rejection Sampling Fine-Tuning (RFT) \citep{yuan2023scaling}, Direct Preference Optimization (DPO) \citep{dpo}, PPO, and GRPO—as direct or streamlined forms of RL. Through this lens, we conduct a comprehensive suite of experiments—such as online versus offline learning, supervision at the outcome versus process level, single-round versus iterative RL, among others—to rigorously explore the key factors within this framework. Finally, we provide an analysis of how our RL methodology enhances the capabilities of instruction-tuned models, and we outline avenues for future research that could lead to further RL advancements within this unified context.

\subsection{Contributions}
We present advancements in large-scale mathematical pre-training and conduct a comprehensive study and evaluation of reinforcement learning methodologies.

\textbf{Math Pre-Training at Scale}

\begin{itemize}[topsep=0pt]
    \item This study demonstrates that the openly available Common Crawl resource harbors significant mathematical data. Through the application of a carefully crafted data selection methodology, we assemble the DeepSeekMath Corpus—a robust, high-quality dataset comprising 120B tokens sourced from math-related web pages. This corpus is nearly sevenfold larger than the dataset employed by Minerva \citep{minerva} and approximately nine times the scale of the newly introduced OpenWebMath \citep{openwebmath}.

\item The DeepSeekMath-Base 7B, our foundational pre-trained model, demonstrates performance on par with Minerva 540B \citep{minerva}, suggesting that sheer parameter count is not the sole determinant of proficiency in mathematical reasoning. Notably, robust results can also be attained by a more compact model trained on superior-quality datasets.

\item We present results from our experiments on mathematics-focused training. Engaging models in code training before mathematics instruction enhances their performance on mathematical problem-solving tasks, regardless of whether external tools are utilized. This provides some evidence addressing the persistent question: \textit{does code training enhance reasoning skills?} Our results indicate it does, particularly when it comes to mathematical reasoning.

\item While the use of arXiv papers for training is widespread, particularly among math-centric studies, such an approach yields no significant enhancement in performance across any of the mathematical benchmarks considered in this work.
\end{itemize}

\textbf{Exploration and Analysis of Reinforcement Learning}

\begin{itemize}[topsep=0pt]
    \item We propose Group Relative Policy Optimization (GRPO), a reinforcement learning algorithm that is both resource-efficient and effective. Unlike Proximal Policy Optimization (PPO), GRPO does not require a critic model; instead, it utilizes group scores to compute the baseline, thereby substantially lowering training overhead.
    \item We show that applying GRPO markedly improves the performance of the instruction-tuned model DeepSeekMath-Instruct using only the instruction-tuning dataset. Additionally, the reinforcement learning process with GRPO yields noticeable improvements on out-of-domain evaluation tasks.
    \item We present a unified framework that contextualizes several reinforcement learning methods, including RFT, DPO, PPO, and GRPO. Our empirical investigations encompass a range of experimental setups—such as comparing online versus offline training, outcome-based versus process-based supervision, and single-step versus iterative reinforcement learning—to thoroughly analyze the critical components within this framework.
    \item Leveraging insights from our unified framework, we delve into the mechanisms driving reinforcement learning’s effectiveness for LLMs, and outline multiple promising avenues for advancing reinforcement learning techniques in this context.
\end{itemize}

\subsection{Summary of Evaluations and Metrics}

\begin{itemize}[topsep=0pt]
    \item \textbf{English and Chinese Mathematical Reasoning}: We systematically evaluate our models across a range of English and Chinese datasets, encompassing mathematical tasks from elementary to collegiate complexity. The English evaluation suite consists of GSM8K \citep{gsm8k}, MATH \citep{MATH}, SAT \citep{llemma}, OCW Courses \citep{minerva}, and MMLU-STEM \citep{mmlu}. For Chinese, the benchmarks considered are MGSM-zh \citep{mgsm}, CMATH \citep{wei2023cmath}, Gaokao-MathCloze \citep{agieval}, and Gaokao-MathQA \citep{agieval}. In our experiments, we measure model performance in two settings: producing fully self-explanatory textual solutions unaided by external tools, and solving mathematical questions by leveraging Python.

On English evaluation tasks, DeepSeekMath-Base achieves performance on par with the proprietary Minerva 540B model \citep{minerva}, while outperforming all currently available open-source base models—such as Mistral 7B \citep{mistral} and Llemma-34B \citep{llemma}—by a considerable margin, independent of the presence of math pre-training. Remarkably, DeepSeekMath-Base demonstrates even greater advantages on Chinese benchmarks. This may be attributed to our decision to go beyond the English-only math pre-training data collection approach found in prior studies \citep{minerva,llemma} and incorporate high-quality math data in other languages as well. When subjected to further mathematical instruction fine-tuning and reinforcement learning, the resulting DeepSeekMath-Instruct and DeepSeekMath-RL models achieve over 50\% accuracy on the competition-level MATH dataset—a milestone reached for the first time in the open-source domain.

\item \textbf{Formal Mathematics}: For the evaluation of DeepSeekMath-Base in the context of informal-to-formal theorem proving, we follow the methodology outlined by \citep{dsp_proof}, using the miniF2F benchmark \citep{minif2f} and the Isabelle proof assistant \citep{isabelle}. DeepSeekMath-Base shows notable effectiveness in few-shot autoformalization scenarios.

\item \textbf{Natural Language Understanding, Reasoning, and Code}: To thoroughly assess the model's broad competencies in language comprehension, logical reasoning, and programming, we employ the Massive Multitask Language Understanding (MMLU) benchmark \citep{mmlu}, which features 57 multiple-choice tasks drawn from a wide spectrum of domains. We also utilize BIG-Bench Hard (BBH) \citep{bbh}, comprised of 23 complex tasks that generally entail multi-step reasoning, along with HumanEval \citep{codex} and MBPP \citep{mbpp}, standard benchmarks for code generation abilities. Pre-training on mathematical data appears to enhance performance not only in language understanding but also in reasoning tasks.

\section{Math Pre-Training}

\subsection{Data Collection and Decontamination} This section details the methodology used to assemble the DeepSeekMath Corpus utilizing Common Crawl data. As shown in Figure \ref{fig:data_collect}, our approach employs an iterative pipeline designed to efficiently extract a large-scale mathematical corpus from Common Crawl, initiated by a high-quality seed corpus composed of select math-focused datasets. Notably, this systematic framework is generalizable and may be extended for use in other domains, including but not limited to coding.

Initially, we select OpenWebMath \citep{openwebmath}, a well-curated repository of mathematical texts from the web, to serve as our seed dataset. Building upon this, we leverage the seed corpus to train a fastText model \citep{joulin2016fasttext} to identify and retrieve additional web pages exhibiting similar mathematical content. For training, we randomly sample 500,000 entries from the seed corpus as positive instances, alongside 500,000 randomly chosen pages from the Common Crawl as negative samples. Training is performed using an open-source fastText implementation\footnote{\url{https://fasttext.cc}}, with the following hyperparameters: vector dimensionality of 256, a learning rate set at 0.1, word n-grams capped at length 3, a minimum threshold of 3 occurrences per word, and three training epochs.

To make the Common Crawl manageable, we utilize both URL-based deduplication and near-duplicate filtering, which condenses the crawl to 40 billion unique HTML pages. The trained fastText model is subsequently applied to the deduplicated crawl to recall mathematical pages. To ensure the quality of mathematical texts retained, we sort the recalled web pages according to their fastText-predicted scores and retain only the highest scoring fraction. We empirically measure how much data to keep by performing pre-training using the top 40B, 80B, 120B, and 160B token subsets. For the first cycle, we decide to preserve the top 40B tokens.

Following the initial round of data acquisition, a substantial number of mathematical web pages remain uncaptured, primarily due to insufficient variety in the positive examples used to train the fastText model. To address this limitation, we supplement the seed corpus by sourcing additional mathematical web domains. The procedure involves partitioning the entirety of Common Crawl into non-overlapping domains, with each domain comprising web pages sharing an identical base URL. For every domain, we compute the fraction of web pages collected during the first round. If this proportion exceeds 10\%, the domain is flagged as math-related (for instance, \url{mathoverflow.net}). Next, we conduct a manual annotation process on the URLs within these recognized domains to pinpoint those specifically associated with mathematical content (such as \url{mathoverflow.net/questions}). Web pages that are linked to these annotated URLs but were previously excluded are then incorporated into the seed corpus. This strategy ensures that we can expand the pool of positive examples, leading to a more robust fastText model with a heightened capability to identify mathematical data in future iterations. After conducting four collection cycles, we accumulate a total of 35.5M mathematical web pages, amounting to 120B tokens. In the fourth round, we observe that almost 98\% of the material had already been collected by the third round, prompting us to terminate the data collection process.

In order to prevent benchmark contamination, we adopt the filtering methodology described in \cite{deepseek-coder}, removing web pages that include content from English mathematical benchmarks like GSM8K \citep{gsm8k} and MATH \citep{MATH}, as well as Chinese datasets such as CMATH \citep{wei2023cmath} and AGIEval \citep{agieval}. Our process eliminates any text fragment that contains an exact 10-gram match with substrings found in the benchmark datasets from the math training set. For instances in which benchmark text spans are shorter than 10-grams but at least 3-grams long, we also utilize exact matching to filter out any web pages compromised by overlap with the evaluation benchmarks.

\subsection{Validating the Quality of the DeepSeekMath~Corpus}

We run pre-training experiments to investigate how the DeepSeekMath~Corpus is compared with the recently released math-training corpora:

\begin{itemize}[topsep=0pt]
    \item \textbf{MathPile} \citep{mathpile}: This dataset is an 8.9B token collection drawn from various sources, including textbooks, Wikipedia, ProofWiki, CommonCrawl, StackExchange, and arXiv, with arXiv accounting for more than 85\% of the tokens;
    \item \textbf{OpenWebMath} \citep{openwebmath}: Comprising 13.6B tokens, this resource consists of CommonCrawl documents that have been selected for mathematical relevance;
    \item \textbf{Proof-Pile-2} \citep{llemma}: This mathematical dataset brings together OpenWebMath, AlgebraicStack (which contains 10.3B tokens of math-focused code), and arXiv papers (with 28.0B tokens). For all experiments involving Proof-Pile-2, we adopt the arXiv:Web:Code distribution ratio of 2:4:1 as recommended in \cite{llemma}.
\end{itemize}

\subsubsection{Training Setting}
\label{sec:quality-policy}
Mathematical training is carried out on a base language model consisting of 1.3B parameters, structured in accordance with the DeepSeek LLMs architecture \citep{deepseek-llm}, and referred to here as DeepSeek-LLM 1.3B. Separate models are fine-tuned on each distinct mathematical dataset, each undergoing 150B token updates. All training procedures are implemented with the streamlined and resource-efficient HAI-LLM framework \citep{haillm}. Consistent with DeepSeek LLMs' methodology, the AdamW optimizer \citep{adamW} is adopted, configured with $\beta_1=0.9$, $\beta_2=0.95$, and $\mathrm{weight\_decay}=0.1$. The learning rate follows a multi-step decay schedule: an initial warmup period of 2,000 steps leads to the maximum learning rate, which subsequently declines to 31.6\% at 80\% of total training and further drops to 10.0\% at 90\% completion. The peak learning rate is capped at 5.3e-4. For each update, a batch containing 4M tokens is processed, with a context window of 4K tokens.

\subsubsection{Evaluation Results}

\textbf{The DeepSeekMath~Corpus is of high quality, covers multilingual mathematical content, and is the largest in size.}

\begin{itemize}[topsep=0pt]
    \item \textbf{High-quality}:  
    To assess the downstream effectiveness, we conduct few-shot chain-of-thought prompting~\cite{cot} across eight mathematical benchmarks. According to Table \ref{tab:corpora_comparison}, models trained using the DeepSeekMath~Corpus exhibit a marked advantage in performance. Moreover, as illustrated by Figure \ref{fig:corpora_comparisons}, after exposure to 50B tokens (the equivalent of a single epoch on Proof-Pile-2), the model trained on DeepSeekMath~Corpus surpasses the counterpart trained on Proof-Pile-2. This result suggests that, on average, the DeepSeekMath~Corpus delivers superior quality data. 
    \item \textbf{Multilingual}:  
    Containing resources in several languages, the DeepSeekMath~Corpus especially emphasizes English and Chinese, which are the predominant languages in its dataset. Table \ref{tab:corpora_comparison} demonstrates that leveraging the DeepSeekMath~Corpus results in improved mathematical reasoning capabilities in both English and Chinese. By comparison, previous math corpora, which are largely centered on English, contribute only marginally, and can even degrade, model performance in Chinese mathematical tasks.
    \item \textbf{Large-scale}:  
    In terms of size, the DeepSeekMath~Corpus substantially exceeds prior mathematical corpora. Figure \ref{fig:corpora_comparisons} reveals that DeepSeek-LLM 1.3B, when trained on DeepSeekMath~Corpus, achieves a more rapid and sustained advancement. Conversely, existing corpora used for baseline models are considerably smaller, and repeated exposure during training causes these models’ performance to plateau much sooner.
\end{itemize}

\subsection{Training and Evaluating DeepSeekMath-Base 7B}

This section presents DeepSeekMath-Base 7B, a foundational model distinguished by its advanced reasoning capacity in mathematical contexts. We initialized the model using DeepSeek-Coder-Base-v1.5 7B \citep{deepseek-coder} and carried out training over a total of 500B tokens. The dataset composition includes 56\% from the DeepSeekMath Corpus, 4\% sourced from AlgebraicStack, 10\% originating from arXiv, 20\% constituted by code from Github, with the final 10\% comprising natural language text extracted from Common Crawl in both English and Chinese. The training procedure largely follows the framework outlined in Section \ref{sec:quality-policy}, with two main adjustments: a peak learning rate of 4.2e-4 is applied, and the token batch size is set at 10M.

We undertake an extensive evaluation of DeepSeekMath-Base 7B's mathematical proficiency, examining its capacity to independently generate complete mathematical solutions, utilize external tools for problem-solving, and perform formal theorem proving tasks. In addition to its mathematical assessment, we offer an overview of the base model's broader capabilities, encompassing natural language understanding, reasoning abilities, and programming competence.

\paragraph{Mathematical Problem Solving with Step-by-Step Reasoning}

We assess DeepSeekMath-Base's capabilities in mathematical problem solving through few-shot chain-of-thought prompting \citep{cot}, utilizing eight distinct benchmarks in both English and Chinese. The selected benchmarks span quantitative reasoning tasks—such as GSM8K \citep{gsm8k}, MATH \citep{MATH}, and CMATH \citep{wei2023cmath}—as well as multiple-choice assessments, including MMLU-STEM \citep{mmlu} and Gaokao-MathQA \citep{agieval}. These benchmarks collectively address a broad spectrum of mathematical disciplines, ranging from elementary to collegiate levels of difficulty.

Table \ref{tab:base_math_cot} highlights that DeepSeekMath-Base 7B achieves the strongest results across all eight evaluated benchmarks when compared to other open-source base models, such as the popular general-purpose Mistral 7B \citep{mistral} and the newly introduced Llemma 34B \citep{llemma}, which is trained on math data from Proof-Pile-2 \citep{llemma}. Remarkably, on the advanced MATH competition dataset, DeepSeekMath-Base exceeds the performance of existing open-source base models by more than 10\% in absolute terms, and even surpasses Minerva 540B \citep{minerva}—a proprietary base model that is 77 times larger, constructed on PaLM \citep{palm} and subjected to further training on mathematical corpora.

\paragraph{Mathematical Problem Solving with Tool Use} We assess the capability for mathematically-focused program-assisted reasoning on the GSM8K and MATH datasets, employing few-shot program-of-thought prompting as introduced by \citet{pot,pal}. The approach requires models to generate solutions for each problem in the form of Python code, leveraging packages like \textit{math} and \textit{sympy} to carry out advanced calculations. The program's executed output is considered the final answer. Table \ref{tab:base_math_pot_proof} indicates that DeepSeekMath-Base 7B achieves superior performance compared to the previous leading model, Llemma 34B.

\paragraph{Formal Mathematics}

Automating formal proofs plays a crucial role in verifying the correctness and dependability of mathematical reasoning, while also improving workflow efficiency—a topic that has recently garnered substantial focus. In our study, we assess DeepSeekMath-Base 7B using the informal-to-formal proving task described in \citep{dsp_proof}, where the objective is to generate a formal proof based on an informal description, its corresponding formalized statement, and an informal proof. Evaluations are performed on the miniF2F benchmark \citep{minif2f}, which targets formalization of Olympiad-level math problems. For each case, we construct a formal Isabelle proof utilizing few-shot prompting. Consistent with the methodology in \cite{dsp_proof}, we prompt models to create proof sketches, then employ the standard automated prover Sledgehammer \citep{sledgehammer} to fill in any remaining proof steps. As indicated in Table \ref{tab:base_math_pot_proof}, DeepSeekMath-Base 7B achieves impressive results in automatic formalization of proofs.

\paragraph{Natural Language Understanding, Reasoning, and Code}

We assess the capabilities of the models in natural language comprehension using MMLU \citep{mmlu}, in reasoning using BBH \citep{bbh}, and in programming through HumanEval \citep{codex} and MBPP \citep{mbpp}. According to Table \ref{tab:understanding_reasoning_code}, DeepSeekMath-Base 7B demonstrates marked improvements over its predecessor, DeepSeek-Coder-Base-v1.5 \citep{deepseek-coder}, particularly in MMLU and BBH, indicating that mathematical training substantially benefits language understanding and reasoning abilities. Furthermore, owing to the continued inclusion of code tokens during training, DeepSeekMath-Base 7B successfully retains the coding benchmark performance levels established by DeepSeek-Coder-Base-v1.5. On the whole, DeepSeekMath-Base 7B achieves considerably better results than the general-purpose model Mistral 7B \citep{mistral} across all three benchmarks pertaining to reasoning and programming.

\section{Supervised Fine-Tuning}

\subsection{SFT Data Curation}

We develop a comprehensive instruction-tuning dataset in mathematics, encompassing both English and Chinese problem statements. These span a wide array of mathematical domains and present a range of difficulty levels. Each problem is matched with a solution articulated in one of several reasoning formats: chain-of-thought (CoT) \citep{cot}, program-of-thought (PoT) \citep{pot,pal}, or a format utilizing external tools \citep{tora}. Altogether, the dataset consists of 776K training samples.

\begin{itemize}[topsep=0pt]
    \item \textbf{English mathematical datasets}: We enhance GSM8K and MATH by providing tool-based annotated solutions, and supplement this with selected problems from MathInstruct \citep{MathInstruct} as well as examples from the Lila-OOD \citep{lila} training partition, wherein solutions utilize CoT or PoT strategies. The English dataset encompasses a broad range of mathematical domains, including but not limited to algebra, probability, number theory, calculus, and geometry.
    \item \textbf{Chinese mathematical datasets}: Our Chinese dataset comprises K-12 mathematics problems that span 76 distinct sub-domains, for example, linear equations; each problem is annotated using both CoT and tool-integrated reasoning approaches.
\end{itemize}

\subsection{Training and Evaluating DeepSeekMath-Instruct 7B}

This section presents DeepSeekMath-Instruct 7B, a model refined through mathematical instruction tuning originating from DeepSeekMath-Base. To create training instances, examples are concatenated at random until the aggregate context reaches 4K tokens. The model is subsequently trained over 500 steps, employing a batch size of 256 and maintaining a fixed learning rate of 5e-5.

We evaluate models' mathematical performance both without and with tool use, on 4 quantitative reasoning benchmarks in English and Chinese. We benchmark our model against the leading models of the time:

\begin{itemize}[topsep=0pt]
    \item \textbf{Closed-source models} consist of: (1) the GPT series, with GPT-4 \citep{gpt4} and GPT-4 Code Interpreter~\footnote{\url{https://openai.com/blog/chatgpt-plugins\#\#code-interpreter}} recognized as the most advanced; (2) Gemini Ultra and Gemini Pro \citep{gemini}; (3) Inflection-2 \citep{inflection-2}; (4) Grok-1~\footnote{\url{https://x.ai/model-card}}; in addition, several recent offerings from Chinese firms, such as (5) Baichuan-3~\footnote{\url{https://www.baichuan-ai.com}} and (6) the new GLM-4~\footnote{\url{https://open.bigmodel.cn/dev/api\#glm-4}}.

    from the GLM family \citep{glm}.

These models are intended for broad application and have typically been subjected to multiple rounds of alignment refinement.
\item \textbf{Open-source models} encompass: general-purpose architectures such as (1) DeepSeek-LLM-Chat 67B \citep{deepseek-llm}, (2) Qwen 72B \citep{qwen}, (3) SeaLLM-v2 7B \citep{seallm}, and (4) ChatGLM3 6B \citep{chatglm3}, along with specialized models designed to advance mathematical capabilities. These include (5) InternLM2-Math 20B~\footnote{\url{https://github.com/InternLM/InternLM-Math}}—an extension of InternLM2 that has undergone focused math pre-training in addition to subsequent instruction tuning, (6) Math-Shepherd-Mistral 7B, which utilizes PPO training \citep{schulman2017proximal} on Mistral 7B \citep{mistral} with the support of a process-supervised reward model, (7) WizardMath models \citep{wizardmath} that strengthen the mathematical reasoning skills of Mistral 7B and Llama-2 70B \citep{llama2} through an evolve-instruct strategy (a variant of instruction tuning with AI-generated instructions) as well as PPO training using datasets such as GSM8K and MATH, (8) MetaMath 70B \citep{metamath}—a Llama-2 70B model fine-tuned on an enriched compilation of GSM8K and MATH, (9) ToRA 34B \cite{tora}, which takes CodeLlama 34B and further tunes it for mathematical reasoning with tool integration, and (10) MAmmoTH 70B \citep{MathInstruct}, created by applying instruction tuning on Llama-2 70B using MathInstruct data.

According to the results in Table \ref{tab:sft_rl_math}, when evaluated without tool usage, DeepSeekMath-Instruct 7B excels in executing multi-step reasoning tasks. On the challenging, competition-level MATH benchmark, our model outperforms all available open-source systems as well as most closed-source alternatives—such as Inflection-2 and Gemini Pro—by a margin of at least 9\% in absolute terms. This superior outcome extends even to competitors with significantly greater model sizes (e.g., Qwen 72B) or those refined via reinforcement learning targeting mathematical reasoning (e.g., WizardMath-v1.1 7B). Although DeepSeekMath-Instruct achieves comparable scores to the proprietary Chinese models GLM-4 and Baichuan-3 on the MATH dataset, its performance still lags behind that of GPT-4 and Gemini Ultra.

When evaluated in an environment that permits the combination of natural language reasoning with programmatic tool use to solve problems, DeepSeekMath-Instruct 7B achieves nearly 60\% accuracy on MATH, outperforming all previously released open-source models. Across additional benchmark datasets, our model performs on par with DeepSeek-LLM-Chat 67B, which was the leading model in this area despite being ten times greater in size.

\subsection{Group Relative Policy Optimization} Following the Supervised Fine-Tuning (SFT) phase, reinforcement learning (RL) has demonstrated significant capability in enhancing the mathematical reasoning skills of LLMs \citep{wang2023math,wizardmath}. In this part, we present Group Relative Policy Optimization (GRPO), a novel and practical RL approach developed to address this objective.

\subsubsection{From PPO to GRPO} Proximal Policy Optimization (PPO) \citep{schulman2017proximal} is an actor-critic RL algorithm that is widely used in the RL fine-tuning stage of LLMs \citep{ouyang2022training}. In particular, it optimizes LLMs by maximizing the following surrogate objective:

\begin{equation}
\footnotesize
    \mathcal{J}_{PPO}(\theta) = \mathbb{E}{[q \sim P(Q), o \sim \pi_{\theta_{old}}(O|q)]} \frac{1}{|o|} \sum_{t=1}^{|o|} \min \left[ \frac{\pi_\theta(o_{t} | q, o_{<t})}{\pi_{\theta_{old}}(o_{t} | q, o_{<t})} A_{t}, \text{clip} \left( \frac{\pi_\theta(o_{t} | q, o_{<t})}{\pi_{\theta_{old}}(o_{t} | q, o_{<t})}, 1 - \varepsilon, 1 + \varepsilon \right)  A_{t} \right] ,
\end{equation}

The above formulation expresses the PPO objective, where the expectation is taken over questions sampled from $P(Q)$ and outputs drawn from the policy $\pi_{\theta_{old}}(O|q)$. Averaging across the length of the output $|o|$, each timestep $t$ involves evaluating the minimum between the policy ratio scaled by the advantage $A_t$ and the advantage multiplied by the clipped policy ratio. The clip operator ensures that the policy update remains within a specified range $[1 - \varepsilon, 1 + \varepsilon]$, thereby controlling the extent of policy deviation during optimization.

Here, $\pi_{\theta}$ denotes the current policy model, while $\pi_{\theta_{old}}$ refers to the previous policy model. The variables $q$ and $o$ correspond to questions and outputs, which are drawn from the question dataset and generated by the old policy $\pi_{\theta_{old}}$, respectively. The parameter $\varepsilon$ serves as a hyper-parameter for clipping, introduced in PPO to enhance training stability. The advantage term $A_t$ is estimated by leveraging Generalized Advantage Estimation (GAE) \citep{gae}, which makes use of future rewards $\{r_{\ge t}\}$ along with a learned value function $V_{\psi}$. Consequently, PPO requires joint training of the value function in addition to the policy model. To counteract excessive optimization of the reward model, it is standard practice to apply a per-token KL penalty relative to a reference model, incorporating this penalty into the reward at each token step, following \citep{ouyang2022training}, as follows:

\begin{equation}
    r_{t} = r_\varphi(q, o_{\le t}) - \beta \log\frac{\pi_{\theta}(o_{t}|q, o_{<t})}{\pi_{ref}(o_{t}|q, o_{<t})},
\label{eq:PPO-reward}
\end{equation}

Here, $r_\varphi$ represents the reward model, $\pi_{ref}$ denotes the reference model—typically the original SFT model—and $\beta$ serves as the scaling factor for the KL divergence penalty.

As the value function employed in PPO is typically another model of comparable size as the policy model, it brings a substantial memory and computational burden. Additionally, during RL training, the value function is treated as a baseline in the calculation of the advantage for variance reduction. While in the LLM context, usually only the last token is assigned a reward score by the reward model, which may complicate the training of a value function that is accurate at each token. To address this, as shown in Figure \ref{fig:grpo}, we propose Group Relative Policy Optimization (GRPO), which obviates the need for additional value function approximation as in PPO, and instead uses the average reward of multiple sampled outputs, produced in response to the same question, as the baseline. More specifically, for each question $q$, GRPO samples a group of outputs $\{o_1, o_2, \cdots, o_G\}$  from the old policy  $\pi_{\theta_{old}}$  and then optimizes the policy model by maximizing the following objective:

\begin{equation}
\footnotesize
\begin{split}
    \mathcal{J}_{GRPO}(\theta) &= \mathbb{E}{[q \sim P(Q), \{o_i\}_{i=1}^G \sim \pi_{\theta_{old}}(O|q)]}  \\
    & \frac{1}{G}\sum_{i=1}^G\frac{1}{|o_i|} \sum_{t=1}^{|o_i|} \left\{ \min \left[ \frac{\pi_\theta(o_{i,t} | q, o_{i,<t})}{\pi_{\theta_{old}}(o_{i,t} | q, o_{i,<t})} \hat{A}_{i,t}, \text{clip} \left( \frac{\pi_\theta(o_{i,t} | q, o_{i,<t})}{\pi_{\theta_{old}}(o_{i,t} | q, o_{i,<t})}, 1 - \varepsilon, 1 + \varepsilon \right)  \hat{A}_{i,t} \right] - \beta \mathbb{D}_{KL}\left[\pi_{\theta} || \pi_{ref}\right]\right\} ,
\end{split}
\label{eq:GRPO-obj}
\end{equation}

where $\varepsilon$ and $\beta$ are hyper-parameters, and $\hat{A}_{i,t}$  is the advantage calculated based on relative rewards of the outputs inside each group only, which will be detailed in the following subsections. The group relative way that GRPO leverages to calculate the advantages, aligns well with the comparative nature of rewards models,  as reward models are typically trained on datasets of comparisons between outputs on the same question. Also note that, instead of adding KL penalty in the reward, GRPO regularizes by directly adding the KL divergence between the trained policy and the reference policy to the loss, avoiding complicating the calculation of  $\hat{A}_{i,t}$. And different from the KL penalty term used in (\ref{eq:PPO-reward}), we estimate the KL divergence with the following unbiased estimator \citep{kl_approx}:

\begin{equation}
\small
    \mathbb{D}_{KL}\left[\pi_{\theta} || \pi_{ref}\right] = \frac{\pi_{ref}(o_{i,t}|q,o_{i,<t})}{\pi_{\theta}(o_{i,t}|q,o_{i,<t})}- \log\frac{\pi_{ref}(o_{i,t}|q,o_{i,<t})}{\pi_{\theta}(o_{i,t}|q,o_{i,<t})} - 1,
\end{equation}

which is guaranteed to be positive.

\begin{algorithm}[t]
  \small
  \caption{Iterative Group Relative Policy Optimization}
  \textbf{Input} Initial policy network $\pi_{\theta_{\text{init}}}$; reward estimators $r_\varphi$; collection of task prompts $\mathcal{D}$;
  hyperparameters $\varepsilon$, $\beta$, $\mu$
  \begin{algorithmic}[1]
    \State Set current policy $\pi_\theta \leftarrow \pi_{\theta_{\text{init}}}$
    \For{iteration = 1, \dots, I}
       \State Set reference policy $\pi_{ref} \leftarrow \pi_{\theta}$
      \For{step = 1, \dots, M}
        \State Select a minibatch $\mathcal{D}_b$ from $\mathcal{D}$
        \State Assign $\pi_{\theta_{old}} \leftarrow \pi_{\theta}$ to preserve the current policy parameters
        \State For each query $q$ in $\mathcal{D}_b$, sample $G$ responses $\{o_i\}_{i=1}^G$ from $\pi_{\theta_{old}} (\cdot \mid q)$ 
        \State For every sampled output $o_i$, evaluate rewards $\{r_i\}_{i=1}^{G}$ via the reward model $r_{\varphi}$
        \State Utilize group relative advantage estimation to determine $\hat{A}_{i,t}$ for the $t$-th token in $o_i$
        \For{GRPO iteration = 1, \dots, $\mu$}
          \State Optimize the policy $\pi_{\theta}$ according to the GRPO criterion as specified in Equation \ref{eq:GC-GRPO}
        \EndFor
      \EndFor
      \State Continuously refine $r_\varphi$ using training with a replay buffer
    \EndFor
  \end{algorithmic}
  \textbf{Output} $\pi_\theta$
  \label{alg:iter-grpo}
\end{algorithm}

\subsubsection{Outcome Supervision RL with GRPO} Specifically, given a question $q$, a set of candidate outputs $\{o_1, o_2, \cdots, o_G\}$ is generated using the previous policy $\pi_{\theta_{old}}$. Each output is evaluated by a reward model, resulting in a set of rewards $\mathbf{r}=\{r_1, r_2, \cdots, r_G\}$ assigned to the respective outputs. These reward values are then standardized within the group by subtracting the mean and dividing by the standard deviation of $\mathbf{r}$. Using outcome supervision, the normalized reward associated with each output $o_i$ is assigned at the sequence's conclusion, with the estimated advantage $\hat{A}_{i, t}$ for every token in that output set to this normalized score: $\hat{A}_{i, t} = \widetilde{r}_i = \frac{r_i- {\rm mean}(\mathbf{r})}{{\rm std}(\mathbf{r})}$. The policy is trained to maximize the objective function described in equation (\ref{eq:GRPO-obj}).

\subsubsection{Process Supervision RL with GRPO}

Outcome supervision typically furnishes a single reward signal at the conclusion of each output sequence, which can be inadequate and inefficient for guiding policy learning in intricate mathematical problem settings. Inspired by \cite{wang2023math}, we further investigate the use of process supervision, where reward feedback is provided at the termination of every intermediate reasoning step. Specifically, for a question $q$ accompanied by $G$ sampled solutions $\{o_1, o_2, \cdots, o_G\}$, a dedicated process reward model evaluates each step in these outputs, assigning rewards as follows: $\mathbf{R} = \{ \{r_1^{index(1)},\cdots,r_1^{index(K_1)}\}, \cdots,  \{r_G^{index(1)},\cdots,r_G^{index(K_G)}\} \}$. Here, $index(j)$ refers to the token position at the conclusion of the $j$-th reasoning step, while $K_i$ indicates the number of such steps for the $i$-th response. To ensure consistency in scaling, we standardize the rewards by subtracting the mean and dividing by the standard deviation across all collected step rewards: $\widetilde{r}_i^{index(j)} = \frac{r_i^{index(j)} - {\rm mean(\mathbf{R})}}{{\rm std(\mathbf{R})}}$.

Afterwards, process supervision computes the token-wise advantage as the cumulative sum of these normalized rewards from each subsequent step, which can be expressed as $\hat{A}_{i, t} = \sum_{index(j) \ge t} \widetilde{r}_i^{index(j)}$. The policy parameters are then updated to maximize the expected return given by the objective in equation (\ref{eq:GRPO-obj}).

\subsubsection{Iterative RL with GRPO} During reinforcement learning, the existing reward model might become inadequate for effectively supervising the updated policy model as training proceeds. To address this, we investigate iterative RL with GRPO. As detailed in Algorithm \ref{alg:iter-grpo}, the iterative GRPO approach involves constructing new reward model training datasets using samples generated from the current policy model. The reward model itself is incrementally updated by incorporating 10\% of past data through a replay buffer. Subsequently, we designate the policy model as the reference model and use the revised reward model to continue training the policy model.

\subsection{Training and Evaluating DeepSeekMath-RL}

Our reinforcement learning experiments are conducted using DeepSeekMath-Instruct 7B as the backbone. The RL training corpus comprises approximately 144,000 chain-of-thought questions associated with GSM8K and MATH datasets, derived from the SFT data. We deliberately omit other SFT questions to focus on assessing RL's effects on benchmark tasks not represented during RL training. Construction of reward model training data follows the methodology outlined in \citep{wang2023math}. Our initial reward model is trained on DeepSeekMath-Base 7B, using a learning rate of $2 \times 10^{-5}$. For policy optimization with GRPO, we set the policy model's learning rate to $1 \times 10^{-6}$ and apply a KL divergence coefficient of $0.04$. Each question receives 64 sampled completions, while both the maximum sequence length and batch size are set to 1024. The policy model undergoes only a single parameter update at the conclusion of each exploration round. Evaluation of DeepSeekMath-RL 7B is performed using the same suite of benchmarks as for DeepSeekMath-Instruct 7B. Specifically, GSM8K and MATH with chain-of-thought prompts are categorized as in-domain evaluations for DeepSeekMath-RL 7B, while all additional benchmarks serve as measures of out-of-domain performance.

As illustrated in Table \ref{tab:sft_rl_math}, the evaluation compares open- and closed-source models engaged in both chain-of-thought and tool-augmented reasoning on English and Chinese test sets. The main findings are as follows: 1) DeepSeekMath-RL 7B achieves 88.2\% accuracy on GSM8K and 51.7\% on MATH when employing chain-of-thought reasoning. These results exceed those achieved by any open-source model in the 7B to 70B parameter range, and outperform most closed-source counterparts as well. 2) Importantly, DeepSeekMath-RL 7B’s training regimen consists solely of chain-of-thought-format instruction tuning data from GSM8K and MATH, beginning with DeepSeekMath-Instruct 7B as a base. Notwithstanding the limited diversity in its training set, DeepSeekMath-RL 7B consistently delivers higher scores than DeepSeekMath-Instruct 7B in all reported metrics, thereby highlighting the advantages introduced by reinforcement learning.

\section{Discussion}
Here, we present and analyze the outcomes from both our pre-training phase and reinforcement learning experiments.

\subsection{Lessons Learnt in Pre-Training}

To begin, we detail our observations during pre-training. Unless explicitly mentioned, all training parameters will follow those described in Section \ref{sec:quality-policy}. It should be emphasized that references to the DeepSeekMath Corpus throughout this section pertain to an 89B-token dataset generated during the second round of data collection.

\subsubsection{Code Training Benefits Mathematical Reasoning}

An influential, though empirically unproven, hypothesis proposes that training on code enhances reasoning capabilities. In this work, we seek to address this idea in part by focusing on mathematics: specifically, we show that code training boosts models' mathematical reasoning skills regardless of whether tools are utilized.

To study how code training affects mathematical reasoning, we experimented with the following two-stage training and one-stage training settings:

\textbf{Two-Stage Training}

\begin{itemize}[topsep=0pt]
    \item \textbf{Code Training for 400B Tokens $\rightarrow$ Math Training for 150B Tokens}: In this approach, DeepSeek-LLM 1.3B is initially trained using 400B code tokens, followed by an additional 150B tokens specifically focused on mathematical data;
    \item \textbf{General Training for 400B Tokens $\rightarrow$ Math Training for 150B Tokens}: For comparative purposes, we conduct a parallel experiment where, in the first phase, the model is exposed to 400B general tokens (sampled from DeepSeek-AI’s extensive general corpus) instead of code tokens, prior to 150B tokens of mathematical content, with the aim of analyzing the relative benefits of code versus general tokens for enhancing mathematical reasoning.
\end{itemize}

\textbf{One-Stage Training}

\begin{itemize}[topsep=0pt]
    \item \textbf{Math Pretraining with 150B Tokens}: DeepSeek-LLM 1.3B undergoes training on a dataset comprising 150 billion mathematics tokens;
    \item \textbf{Simultaneous Training on 400B Code Tokens and 150B Math Tokens}: Sequential math finetuning after code training negatively affects coding abilities. We explore if integrating code and math tokens in a unified training stage can enhance both mathematical reasoning and help prevent catastrophic forgetting.
\end{itemize}

\paragraph{Results}
The downstream outcomes across various training configurations are presented in Table \ref{tab:code-math} and Table \ref{tab:code-math-general-eval}.

Code training improves program-aided mathematical reasoning performance across both two-stage and one-stage training paradigms. As demonstrated in Table \ref{tab:code-math}, engaging solely in code training during the first stage already substantially boosts the capability to address GSM8K and MATH tasks with Python. Introducing math training in the subsequent stage brings about additional progress. Notably, in the one-stage training scenario, jointly incorporating code and math tokens not only prevents catastrophic forgetting often observed with two-stage training, but also enhances performance in both pure coding tasks (Table \ref{tab:code-math-general-eval}) and program-aided mathematical reasoning tasks (Table \ref{tab:code-math}).

Training with code also enhances mathematical reasoning skills even in the absence of tool use. Within the two-stage training framework, the early phase of code training alone yields notable improvements. This approach further accelerates and amplifies the effectiveness of later math-specific training, culminating in superior overall results. Conversely, merging code tokens and math tokens in a single-stage training process appears to hinder the development of mathematical reasoning abilities when not using tools. One possible explanation is that DeepSeek-LLM 1.3B, owing to its relatively small size, may not possess sufficient capacity to thoroughly integrate both code-related and mathematical knowledge during a unified training phase.

\subsubsection{ArXiv Papers Seem Ineffective in Improving Mathematical Reasoning}

ArXiv papers are commonly included as a component of math pre-training data \citep{minerva,gpt-f,llemma,mathpile}. However, detailed analysis regarding their impact on mathematical reasoning has not been extensively conducted. Perhaps counter-intuitively, according to our experiments, arXiv papers seem ineffective in improving mathematical reasoning. We experiment with models of different sizes, including DeepSeek-LLM 1.3B and DeepSeek-Coder-Base-v1.5 7B \citep{deepseek-coder}, using arXiv corpora that underwent varied processing pipelines:

\begin{itemize}[topsep=0pt]
    \item \textbf{MathPile} \citep{mathpile}: comprises 8.9B tokens curated via heuristic-based cleaning and filtering processes, with more than 85\% of its content sourced from scientific papers on arXiv;
    \item \textbf{ArXiv-RedPajama} \citep{redpajama}: consists of the full collection of arXiv LaTeX documents, excluding preambles, comments, macros, and bibliographies, resulting in a dataset of 28.0B tokens.
\end{itemize}

For our experimental setup, DeepSeek-LLM 1.3B was individually trained for 150B tokens, and DeepSeek-Coder-Base-v1.5 7B for 40B tokens, utilizing each arXiv-specific dataset. Our findings suggest that arXiv papers do not contribute effectively to enhancing mathematical reasoning. Training exclusively on arXiv data failed to yield any substantial gains and, in some cases, resulted in diminished performance across a range of mathematical evaluation benchmarks considered in this work. These benchmarks span quantitative reasoning datasets such as GSM8K and MATH (Table \ref{tab:arxiv-cot}), various multiple-choice assessments exemplified by MMLU-STEM (Table \ref{tab:arxiv-cot}), and formal mathematics benchmarks like miniF2F (Table \ref{tab:arxiv_atp}).

However, this conclusion has its limitations and should be taken with a grain of salt. We have not yet studied:

\begin{itemize}[topsep=0pt]
    \item The influence of arXiv tokens on particular mathematics-oriented tasks that fall outside the scope of this study, for example, the process of transforming formal theorems or proofs into their informal representations;
    \item The consequences of integrating arXiv tokens with additional data sources;
    \item The extent to which the advantages provided by arXiv papers become evident when utilized in models of greater scale.
\end{itemize}

Thus, further exploration is required, which we leave for future studies.

\subsection{Insights of Reinforcement Learning}
\subsubsection{Towards to a Unified Paradigm} Here, we introduce a comprehensive framework for examining various training strategies, including SFT, RFT, DPO, PPO, and GRPO. Additionally, we perform experimental investigations to better understand the elements contributing to this unified framework. Broadly speaking, for a given training algorithm, the gradient with respect to its parameter $\theta$ can be expressed as:

\begin{equation}
    \nabla_{\theta}\mathcal{J}_{\textcolor{red}{\mathcal{A}}}(\theta) = \mathbb{E}[\underbrace{(q,o) \sim \textcolor{red}{\mathcal{D}}}_{Data \ Source}]\left( \frac{1}{|o|} \sum_{t=1}^{|o|}  \underbrace{GC_{{\mathcal{A}}}(q, o, t, \textcolor{red}{\pi_{{rf}}})}_{Gradient \ Coefficient}  \nabla_{\theta}\log \pi_{\theta}(o_t | q, o_{<t})\right).
\label{eq:objective}
\end{equation}

There exist three key components: 1) \textit{Data Source $\mathcal{D}$}, which determines the training data; 2) \textit{Reward Function $\pi_{{rf}}$}, which is the source of the training reward signal; 3) \textit{Algorithm $\mathcal{A}$}: which processes the training data and the reward signal to the gradient coefficient $GC$ that determines the magnitude of the penalty or reinforcement for the data. We analyze several representative methods based on such a unified paradigm:

\begin{itemize}
    \item \textbf{Supervised Fine-tuning (SFT)}: This method involves additional training of a pretrained model utilizing SFT datasets curated by humans.
    \item \textbf{Rejection Sampling Fine-tuning (RFT)}: RFT extends the SFT procedure by conducting further training on a subset of outputs produced by the SFT model in response to SFT questions, retaining only those outputs whose answers are validated as correct.
    \item \textbf{Direct Preference Optimization (DPO)}: DPO further adapts the SFT model by fine-tuning it on supplementary outputs, generated by the SFT model, and employs a pairwise DPO loss for optimization.
    \item \textbf{Online Rejection Sampling Fine-tuning (Online RFT)}: In contrast to standard RFT, Online RFT uses the SFT model to initialize the policy and continually fine-tunes the policy on outputs sampled dynamically from the current policy model.
    \item \textbf{PPO/GRPO}: PPO/GRPO methods also initialize the policy with the SFT model and enhance it by reinforcing outputs drawn from the evolving real-time policy model.
\end{itemize}

Table \ref{tab:data_policy} provides an overview of the elements involved in these methods. A comprehensive explanation of the derivation process is available in Appendix \ref{app:analysis-rl}.

\paragraph{Observation about Data Source} The data source can be categorized into two types: online sampling and offline sampling. Online sampling refers to training data acquired from the live exploratory outcomes generated by the policy model during training, whereas offline sampling indicates that training data originates from samples collected using the initial SFT model. Both RFT and DPO utilize offline sampling approaches, in contrast to Online RFT and GRPO, which employ online sampling methods.

As illustrated in Figure \ref{fig:rl-analysis}, our results indicate that Online RFT achieves markedly better performance than RFT across both benchmarks. In the initial training phases, Online RFT and RFT perform similarly; however, as training proceeds, Online RFT establishes a clear lead, underscoring the benefits of online training. This outcome aligns with expectations: at the start, the actor and the SFT model are nearly identical, so the sampled data reflect only slight variations. Over time, as the actor diverges more from the SFT model, the sampled data begin to display larger differences, at which point the advantages of real-time data sampling with Online RFT become pronounced.

\paragraph{Observation about Gradient Coefficient} The algorithm utilizes the reward signal to determine the gradient coefficient, which subsequently influences the update of model parameters. In our experiments, we separate the reward function into two distinct categories: `Rule' and `Model'. Here, `Rule' entails evaluating the response quality solely by assessing answer correctness, whereas `Model' involves employing a reward model that assigns a score to each response. Notably, the reward model's training data originates from rule-based assessments. As demonstrated in Equations \ref{eq:GC-RFT} and \ref{eq:GC-GRPO}, there exists a fundamental distinction between GRPO and Online RFT. GRPO stands out by dynamically tuning its gradient coefficient according to the reward values generated by the reward model, thus enabling tailored positive or negative reinforcement depending on the reward magnitude of each response. By contrast, Online RFT does not possess this adaptability; it refrains from penalizing incorrect responses and consistently applies an equal degree of positive reinforcement to all responses that yield correct answers.

As shown in Figure \ref{fig:rl-analysis}, GRPO outperforms online RFT, emphasizing the effectiveness of adjusting gradient coefficients for positive and negative cases. Moreover, the results for GRPO+PS exceed those for GRPO+OS, demonstrating the advantages of employing more nuanced, step-specific gradient coefficients. Additionally, we investigate the impact of iterative RL by running two successive iterations in our experiments. Figure \ref{fig:iter-rl} illustrates that iterative RL leads to notable performance gains, particularly during the initial iteration.

\subsubsection{Why RL Works?} In this work, we apply reinforcement learning to a selected portion of instruction tuning data, which substantially improves upon the baseline performance of the instruction-tuned model. To further elucidate the mechanisms underlying this improvement, we compare the Pass@K and Maj@K accuracies of both the Instruct model and its RL-optimized counterpart across two benchmark datasets. Figure \ref{fig:maj-pass} demonstrates that while RL markedly enhances Maj@K accuracy, there is little to no effect on Pass@K. These results suggest that the principal contribution of RL lies in strengthening the overall output distribution—specifically, \textbf{the gain can largely be traced to an increased likelihood of producing the correct answer within the TopK outputs, rather than an intrinsic advancement in base capabilities.} In a similar vein, \citep{wang2023making} point to a \textbf{misalignment problem} affecting reasoning tasks in SFT models, and show that these models' reasoning skills are improvable via preference alignment strategies \citep{yuan2023rrhf,song2023preference,wang2023making}.

\subsubsection{How to Achieve More Effective RL?}
\label{sec:effec_rl}
Our findings indicate that RL performs effectively on mathematical reasoning challenges. Additionally, we introduce a cohesive framework that accommodates various well-known training strategies, framing them as instances of either standard or reduced RL approaches. As captured in Equation \ref{eq:objective}, three fundamental elements are involved: Data Source, Algorithm, and Reward Function. We discuss possible avenues for advancing each of these components in future research.

\paragraph{Data Source} The data source serves as the foundational input for all training techniques. In the realm of RL, the data source specifically denotes the collection of unlabeled questions along with corresponding outputs, which are generated through the policy model. For this study, our approach involves exclusively utilizing questions derived from the instruction tuning phase and employing a straightforward nucleus sampling approach for generating responses. We hypothesize that this may be a contributing factor as to why our RL methodology yields improvements predominantly in Maj@K performance. Looking ahead, we intend to extend our RL pipeline to handle out-of-distribution question prompts and to incorporate \textbf{advanced sampling (decoding) strategies}, such as those leveraging tree-search methodologies \citep{yao2023tree}. Furthermore, advancements in \textbf{efficient inference techniques} \citep{xia-etal-2023-speculative,leviathan2023fast,kwon2023efficient,xia2024unlocking}, which govern the efficiency with which policy models explore the solution space, are also of critical significance.

\paragraph{Algorithms} The algorithms leverage the data and reward signals to calculate gradient coefficients, facilitating updates to the model parameters. As delineated in Equation \ref{eq:objective}, most contemporary approaches entirely \textbf{TRUST} the reward function’s signal when adjusting the conditional probability for specific tokens. Nonetheless, in practice, the reward signal’s dependability cannot always be guaranteed, particularly in highly intricate tasks. For instance, the PRM800K datasets \citep{lightman2023let}, despite being meticulously annotated by experienced annotators, still exhibit an error rate of roughly 20\% in their labels\footnote{\url{https://github.com/openai/prm800k/issues/12\#issuecomment-1728491852}}. Given these challenges, our focus shifts to developing reinforcement learning algorithms that maintain robustness even under noisy reward conditions. We contend that \textbf{WEAK-TO-STRONG} \citep{burns2023weak} alignment strategies have the potential to fundamentally reshape learning algorithm design.

\paragraph{Reward Function} The reward function serves as the primary source of feedback during training. In reinforcement learning (RL), this typically takes the form of a neural reward model. There are three significant avenues for advancing reward models: 1) \textbf{Enhancing the generalization capabilities of reward models.} For reinforcement learning to bolster the core capacities of large language models (LLMs)—rather than just preserving their output distributions—the reward model must be robust enough to address unseen questions and novel decoding results; 2) \textbf{Capturing and representing reward model uncertainty.} Modeling uncertainty can serve as an important connection between unreliable reward signals and learning approaches designed to transfer from weaker to stronger models; 3) \textbf{Constructing high-quality process reward models in a data-efficient manner} that supply rich, step-wise feedback to support complex reasoning tasks \citep{lightman2023let,wang2023math}.

\section{Conclusion, Limitation, and Future Work} This paper introduces DeepSeekMath, an open-source model that surpasses existing open models on the competition-level MATH benchmark and nearly matches the results obtained by proprietary systems. DeepSeekMath~leverages a DeepSeek-Coder-v1.5 7B foundation and is continually trained with 500B tokens, including a substantial portion—120B math-focused tokens—extracted from Common Crawl. Through detailed ablation analyses, we find that data from web pages provides a considerable resource for acquiring quality mathematical content, whereas contributions from arXiv are less substantial than anticipated. Additionally, we propose Group Relative Policy Optimization (GRPO), which builds upon Proximal Policy Optimization (PPO) and delivers notable gains in mathematical reasoning with reduced memory requirements. Experimental evaluation demonstrates that GRPO remains impactful even when DeepSeekMath-Instruct 7B already performs strongly across benchmarks. Finally, we outline a unified conceptual framework for interpreting related methods and suggest various promising avenues for enhancing reinforcement learning strategies.

While DeepSeekMath~demonstrates strong performance on benchmarks focused on quantitative reasoning, its proficiency in geometry and theorem-proving tasks remains below that of closed-source models. As observed during preliminary testing, the model struggles with problems involving geometric concepts such as triangles and ellipses, potentially reflecting a bias in the selection of training data during both pre-training and fine-tuning. Furthermore, due to limitations in model size, DeepSeekMath~lags behind GPT-4 in terms of few-shot learning; whereas GPT-4 benefits from the inclusion of few-shot examples, DeepSeekMath~exhibits comparable outcomes whether evaluated under zero-shot or few-shot settings. Moving forward, we intend to enhance the data selection mechanisms to build a more robust and diverse pre-training dataset. Additionally, we plan to investigate the possibilities for more efficient reinforcement learning strategies for large language models, as outlined in Section \ref{sec:effec_rl}.

\end{document}